\documentclass[10pt,a4paper]{article}
\usepackage[T1]{fontenc}
\usepackage[utf8]{inputenc}
\usepackage{lmodern,amsmath,booktabs,graphicx,geometry,caption,titlesec,xcolor,microtype,hyperref}
\geometry{margin=22mm,top=19mm,bottom=19mm}
\hypersetup{hidelinks,pdftitle={What Do Panel Estimators Recover?},pdfauthor={Kelvin Schouten}}
\setlength{\parindent}{0pt}
\setlength{\parskip}{3.5pt}
\definecolor{sectionlabel}{RGB}{99,106,113}
\definecolor{sectionrule}{RGB}{207,211,214}
\renewcommand{\thesection}{\ifnum\value{section}<10 0\fi\arabic{section}}
\titleformat{\section}[display]
  {\normalfont}
  {\color{sectionlabel}\sffamily\fontsize{8}{10}\selectfont SECTION \thesection}
  {2pt}
  {\fontsize{12.8}{15}\selectfont}
  [\vspace{3pt}\color{sectionrule}\titlerule]
\titlespacing*{\section}{0pt}{9pt}{6pt}
\captionsetup{font=small,labelfont=normalfont}
\renewcommand{\arraystretch}{1.12}
\pagestyle{plain}
\begin{document}
\begin{center}
{\Large What Do Panel Estimators Recover?\par}
\vspace{5pt}
{\normalsize A small Monte Carlo experiment with individual effects and persistent errors\par}
\end{center}
\vspace{4pt}

\section{The idea}
I wanted to see how much difference the choice of panel-data estimator makes when a regressor is related to an unobserved individual effect. Rather than use a real dataset, I generated a simple panel where I know the true coefficient. This makes it easy to see which estimators are getting close to the right value and which ones are not.

For individual $i$ in period $t$, the data are generated by
\begin{align}
 y_{it} &= 1+\alpha_i+2x_{it}+u_{it},\\
 x_{it} &= v_{it}+0.5\alpha_i,\\
 u_{it} &= u_{i,t-1}+e_{it}, \qquad u_{i1}=e_{i1}.
\end{align}
Here $\alpha_i$, $v_{it}$, and $e_{it}$ are independent standard normal draws, apart from the dependence built into $x_{it}$ and $u_{it}$. The individual effect $\alpha_i$ stays constant over time, while the error $u_{it}$ follows a random walk. I use $N=200$ individuals, $T=5$ periods, and 1,000 simulated datasets. The true slope is $2$.

\section{Pooled OLS, fixed effects, and first differences}
Pooled OLS ignores $\alpha_i$. Since $x_{it}$ contains part of that effect, the estimated coefficient is biased upwards. In fact, its large-sample target is
\begin{equation}
 2+\frac{\operatorname{Cov}(x_{it},\alpha_i)}{\operatorname{Var}(x_{it})}
 =2+\frac{0.5}{1.25}=2.4.
\end{equation}
Fixed effects (FE) takes away the individual mean, while first differences (FD) subtracts the previous period. Both remove the constant individual effect. In this example FD has another advantage: $\Delta u_{it}=e_{it}$, so differencing also removes the persistence in the error.

\begin{table}[h]
\centering
\caption{Results across 1,000 simulated panels; true coefficient $\beta=2$.}
\begin{tabular}{lrrr}
\toprule
Estimator & Mean & Standard deviation & Bias \\
\midrule
Pooled OLS & 2.392 & 0.074 & +0.392 \\
Fixed effects & 1.997 & 0.035 & $-0.003$ \\
First differences & 1.998 & 0.025 & $-0.002$ \\
\bottomrule
\end{tabular}
\end{table}

The pattern is fairly clear. Pooled OLS is concentrated around the wrong number, while FE and FD are centred close to $2$. FD also varies less across simulations, which is what I would expect with these random-walk errors. This does not mean that FD is always better; the result depends on how the errors were generated.

\begin{figure}[!h]
\centering
\includegraphics[width=0.55\linewidth]{panel_estimates.png}
\caption{Estimated slopes over 1,000 simulations. The dashed line is the true slope, $2$.}
\end{figure}

\newpage
\section{What if the outcome also depends on its past?}
I also tried a dynamic version of the model by adding the previous value of $y$:
\begin{equation}
 y_{it}=1+\alpha_i+2x_{it}+0.5y_{i,t-1}+u_{it}, \qquad t\geq2.
\end{equation}
For the first period I use the original model without the lag. I then estimate the coefficient on $y_{i,t-1}$ using FE and repeat the experiment for different numbers of periods. Each setting again uses 1,000 independently generated panels.

\begin{table}[h]
\centering
\caption{Mean FE estimate of the lag coefficient; true value $\rho=0.5$.}
\begin{tabular}{rr}
\toprule
Periods ($T$) & Mean estimated $\rho$ \\
\midrule
5 & 0.574 \\
10 & 0.691 \\
20 & 0.806 \\
50 & 0.910 \\
\bottomrule
\end{tabular}
\end{table}

\begin{figure}[!h]
\centering
\includegraphics[width=0.55\linewidth]{panel_dynamic.png}
\caption{Average FE estimate of the lag coefficient as $T$ increases.}
\end{figure}

This was the more surprising result for me. The estimate moves further above $0.5$ as I add periods, even though it becomes less variable across simulations. FE removes the fixed individual effect, but the lagged outcome is still related to the error term. The random-walk errors make this different from the usual dynamic-panel example with serially uncorrelated disturbances. So this upward pattern is a feature of this particular simulation, not a general result for FE.

\section{What I take from it}
The main lesson is that a precise estimate does not have to be a correct one. With pooled OLS, more data would make the estimate settle around $2.4$ rather than $2$. FE and FD work much better here because the regressor shocks are independent of the error process once the individual effect is removed. Robust or clustered standard errors can help with inference, but cannot fix the omitted-variable bias in pooled OLS.

I kept the setup intentionally small. The conclusions are about this artificial model, not about panel estimators in every application. The accompanying Python script generates the data, estimates the coefficients, and saves both plots using NumPy and Matplotlib.

{\small The initial data-generating process was used in a coursework assignment completed with Walid Afasi (2026). I independently rewrote and reran the simulations for this short note. For general background, see Wooldridge, \emph{Econometric Analysis of Cross Section and Panel Data} (2010), and Nickell, "Biases in Dynamic Models with Fixed Effects" (1981).}
\end{document}
